\section{And Tell the Router}
\label{sec:ch12-and-tell-the-router}

\index{router!the graph after the guard lands|(}%
Deploy section~\ref{sec:ch12-guard-the-data} and the graph stops working for
everybody. Ask the router for a session and its speaker's address, with a
token the router itself verified:

\begin{minted}[fontsize=\footnotesize,breakanywhere]{graphql}
{ sessions(first: 10) { nodes { title speaker { name email } } } }
\end{minted}

\begin{minted}[fontsize=\footnotesize,breakanywhere]{json}
{"errors":[
  {"message":"Failed to fetch from Subgraph 'speakers' at Path 'sessions.nodes.@.speaker'.",
   "extensions":{
     "errors":[
       {"message":"The current user is not authorized to access this resource.",
        "path":["sessions","nodes","@","speaker","email"],
        "extensions":{"code":"AUTH_NOT_AUTHENTICATED"}},
       {"message":"The current user is not authorized to access this resource.",
        "path":["sessions","nodes","@","speaker","email"],
        "extensions":{"code":"AUTH_NOT_AUTHENTICATED"}},
       {"message":"The current user is not authorized to access this resource.",
        "path":["sessions","nodes","@","speaker","email"],
        "extensions":{"code":"AUTH_NOT_AUTHENTICATED"}}],
     "serviceName":"speakers",
     "statusCode":200}},
  {"message":"Cannot return null for non-nullable field 'Query.sessions.nodes.speaker.name'.",
   "path":["sessions","nodes",0,"speaker","name"]},
  {"message":"Cannot return null for non-nullable field 'Query.sessions.nodes.speaker.name'.",
   "path":["sessions","nodes",1,"speaker","name"]},
  {"message":"Cannot return null for non-nullable field 'Query.sessions.nodes.speaker.name'.",
   "path":["sessions","nodes",2,"speaker","name"]},
  {"message":"Cannot return null for non-nullable field 'Query.sessions.nodes.speaker.name'.",
   "path":["sessions","nodes",3,"speaker","name"]}],
 "data":{"sessions":{"nodes":[
   {"title":"Schemas That Outlive Their Authors","speaker":null},
   {"title":"Reading a Query Plan","speaker":null},
   {"title":"The Bank That Split Its Graph","speaker":null},
   {"title":"Paging Without Offsets","speaker":null}]}}}
\end{minted}

\index{router!three refusals for four sessions}%
Three refusals from the subgraph and four from the router, for four sessions.
The subgraph counts distinct speakers, because
chapter~\ref{ch:router-n-plus-1} established that the router de-duplicates
representations before it sends them and this page has three speakers across
its four sessions. The non-null errors count sessions, because that is where
the nulls land.

\index{router!an anonymous caller is told the same thing}%
Send the same query with no token at all and the answer is byte for byte
identical. That is the tell, and it is the reason this is worth a section:
a graph whose authorization is working looks different to those two callers.
This one does not, because neither of them is getting through and both for
the same reason.

\index{router!headers are not forwarded by default}%
The router does not forward request headers to subgraphs unless it is told
to, and WunderGraph says so in as many words: \enquote{by default, no headers
are forwarded for security reasons}~\autocite{cosmoheaders}. The Speakers
service is not refusing your caller. It is refusing an anonymous request,
because an anonymous request is what reached it.

\index{non-null!the guard's blast radius through the seam}%
The four errors underneath are what a non-null field costs when the entity
carrying it disappears. Losing \code{email} loses the whole \code{Speaker},
and \code{name} is \code{String!}, so every session on the page reports a null
it cannot hold. The titles arrive; nothing under them does.

\index{router!what the missing block survives}%
The part that makes this expensive to diagnose is how narrow it is. The
refusal follows the guarded column rather than the seam, so the same seam
asked for a speaker's public name never selects \code{email}, nothing objects,
and the answer is complete. Every dashboard stays green and one screen is
broken.

\subsection{Telling the Router Who You Are}

\index{router!validating a token}%
The second thing the router needs is a way to know a token at all, and it
takes either a JWKS url, which is what a deployment with an identity provider
behind it uses, or a shared secret. Here is \code{config.yaml} again in full,
the file chapter~\ref{ch:enter-the-router} wrote, with everything this chapter
adds to it marked:

\begin{minted}[highlightlines={19-48}]{yaml}
version: "1"

# Without this the router asks Cosmo Cloud for its execution config and
# refuses to start, because it has no token to ask with. The path is
# resolved against the directory the router is started in, not against
# this file, so start it from this folder:
#
#   router -config config.yaml
execution_config:
  file:
    path: ../graph/router.json

# A development switch, and what makes the query plan visible at all.
# The plan is an Advanced Request Tracing option, and outside dev mode
# the router wants a Cosmo Cloud token before it hands one out. With
# this off, X-WG-Include-Query-Plan is ignored rather than refused.
dev_mode: true

# Who the router thinks is asking. The alternative entry in this list
# is a jwks url, which is what a deployment with an identity provider
# behind it uses; the shared secret is here because this book stands
# up no such provider. header_key_id is not optional: the router
# matches it against the kid in the token header, so the token this
# book mints carries one.
#
# A request carrying a token this key does not verify is answered 401
# by the router and never reaches a subgraph.
authentication:
  jwt:
    jwks:
      - secret: splitting-the-graph-development-signing-key-32b
        symmetric_algorithm: HS256
        header_key_id: dev

# The router does not forward request headers to subgraphs unless it
# is told to, and a subgraph that guards its own data cannot tell who
# is asking without this block. Leaving it out does not make the graph
# insecure; it makes the guarded field unreachable. Every request that
# names one comes back null with an authorization error underneath,
# for a caller holding a good token exactly as for a caller holding
# none, which is what makes the missing block read like a policy
# rather than like a header. A request that never names a guarded
# field is unaffected, so this survives a smoke test.
headers:
  all:
    request:
      - op: propagate
        named: Authorization
\end{minted}

\index{router!a token it cannot verify never reaches a subgraph}%
With that in place the router refuses a token it cannot verify itself, and it
is the first time in this book that the router refuses anything outright
rather than answering 200 with the problem inside the envelope:

\begin{minted}[fontsize=\footnotesize,breakanywhere]{json}
HTTP/1.1 401 Unauthorized

{"errors":[{"message":"unauthorized"}]}
\end{minted}

\index{router!the secret is written twice}%
The secret is in two files, and there is no way around that. The service
validates it in C\# and the router validates it in YAML, and nothing is
shared between a service and a config file. In a deployment both read it from
the same secret store; in this book both have it written out, so that you can
type the system from the page and mint a token that works against either.

\subsection{The Directive, and What It Actually Does}

\index{authenticated@\code{"@authenticated}!the second attribute}%
That leaves the second attribute,
\code{[Authenticated]}, which section~\ref{sec:ch12-guard-the-data} put on
the column and did not explain. It is Apollo's directive rather than Hot
Chocolate's, and the point of it is that it survives composition. Hot
Chocolate emits it and adds it to the import list without being asked:

\begin{graphqlsdl}[fontsize=\footnotesize]
@link(
  url: "https://specs.apollo.dev/federation/v2.6"
  import: ["@authenticated", "@key", "@shareable", "@tag", "FieldSet"]
)
\end{graphqlsdl}

\begin{graphqlsdl}[fontsize=\footnotesize]
email: String! @authorize @authenticated
\end{graphqlsdl}

\index{authenticated@\code{"@authenticated}!what the composer does with it}%
\code{wgc router compose} accepts both, which is worth saying out loud after
chapter~\ref{ch:composition} lost a whole section to a directive Hot
Chocolate emitted and the composer would not take. \code{@authorize} rides
along untouched. \code{@authenticated} is compiled into a rule the router
consumes:

\begin{minted}[fontsize=\footnotesize,breakanywhere]{json}
{"typeName":"Speaker","fieldName":"email",
 "authorizationConfiguration":{"requiresAuthentication":true}}
\end{minted}

\index{authenticated@\code{"@authenticated}!on a type it compiles to the edges}%
Put the directive on a type instead of a field and the shape of that output
does not follow the shape of what you wrote. The composer emits one rule per
field that \emph{returns} the type. So the directive on \code{Speaker}
produces three rules, on \code{Session.speaker}, on the connection's
\code{nodes} and on the edge's \code{node}, and not one of them mentions
\code{Speaker}. Worth knowing before you go looking for your type in that
file.

\subsection{It Is a Filter, Not a Gate}

\index{authenticated@\code{"@authenticated}!the router refuses in its own words}%
The router does enforce the rule. Take Hot Chocolate's attribute back off the
column, leave Apollo's on it, and an unauthenticated caller through the
router is refused:

\begin{minted}[fontsize=\footnotesize,breakanywhere]{json}
{"errors":[
  {"message":"Unauthorized to load field 'Query.sessions.nodes.speaker.email', Reason: not authenticated.",
   "path":["sessions","nodes",0,"speaker","email"],
   "extensions":{"code":"UNAUTHORIZED_FIELD_OR_TYPE"}},
  {"message":"Unauthorized to load field 'Query.sessions.nodes.speaker.email', Reason: not authenticated.",
   "path":["sessions","nodes",1,"speaker","email"],
   "extensions":{"code":"UNAUTHORIZED_FIELD_OR_TYPE"}},
  {"message":"Unauthorized to load field 'Query.sessions.nodes.speaker.email', Reason: not authenticated.",
   "path":["sessions","nodes",2,"speaker","email"],
   "extensions":{"code":"UNAUTHORIZED_FIELD_OR_TYPE"}},
  {"message":"Unauthorized to load field 'Query.sessions.nodes.speaker.email', Reason: not authenticated.",
   "path":["sessions","nodes",3,"speaker","email"],
   "extensions":{"code":"UNAUTHORIZED_FIELD_OR_TYPE"}}],
 "data":{"sessions":{"nodes":[
   {"title":"Schemas That Outlive Their Authors","speaker":null},
   {"title":"Reading a Query Plan","speaker":null},
   {"title":"The Bank That Split Its Graph","speaker":null},
   {"title":"Paging Without Offsets","speaker":null}]}}}
\end{minted}

\index{authenticated@\code{"@authenticated}!no subgraph objected}%
That refusal is the router's own, in its own words, and no subgraph error is
relayed alongside it, because no subgraph objected to anything. The service
has no rule on that column any more.

\index{statement count!the router reads the row and discards it}%
Which raises the question the rest of this book has trained you to ask. What
did that request cost? One statement in the Speakers service. The
entity fetch went out, the reference resolver ran, SQLite returned the row
with the address on it, and the router took the value out of the answer on
the way to the client.

\index{authenticated@\code{"@authenticated}!documented as a post-filter}%
That is not a defect and WunderGraph documents it plainly. Fields carrying
the directive, it says, \enquote{are authorized after the subgraph fetch and
filtered out of the response}~\autocite{cosmoauthentication}. Apollo describes
the same shape for its own router, saying of a protected key field that
\enquote{the router would use the \code{id} field to resolve the
\code{Product} entity, but it wouldn't return
it}~\autocite{apolloauthorization}. Both gateways read first and decide
second.

\index{authenticated@\code{"@authenticated}!what a post-filter cannot protect}%
So what the directive protects is the response. It does not protect the row,
which was read; it does not protect the wire between the router and the
service, which carried the address; and it does not protect port 5002 at all.
With only the directive in place, the entity route on that service hands the
column to anyone who can reach it, with a token or without one, because the
only rule in the graph lives in a process the request never goes through.

\index{authorization!the switch that looks like the answer}%
The router has an authorization option, off by default, whose documentation
says it authorizes an operation's protected fields before any subgraph fetch
runs. That is the description of a gate, and it is the one thing that would
change the paragraph above. Setting it did not change the statement count.
This book therefore ships without it, and nothing in this chapter rests on
it; the switch is named in the research note rather than here, because a
reader who goes and sets it will get the graph I measured rather than the
graph its description promises.

\index{authorization!both, in both places}%
Ship both attributes. \code{[Authenticated]} moves the refusal to the edge
for the ordinary case, gives the client one consistent error whichever
subgraph owns the field, and is the only one of the two that a router can act
on. \code{[Authorize]} is the one that means it. If you have to drop one,
drop the one in the supergraph, because the other is the one still standing
when somebody reaches the port.
\index{router!the graph after the guard lands|)}%
